\documentclass[12pt]{article}

%==================== PACKAGES ====================
\usepackage[a4paper,margin=0.8in]{geometry}
\usepackage{amsmath,amssymb,amsthm}
\usepackage{enumitem}
\usepackage{bm}
\usepackage{hyperref}
\usepackage{fancyhdr}
\usepackage{xcolor}
\usepackage{array}
\usepackage{booktabs}
\usepackage{tikz}
\usepackage{pgfplots}
\pgfplotsset{compat=1.17}
\usepackage[most]{tcolorbox}
\tcbuselibrary{theorems,skins,breakable}
\usepackage{graphicx}
\usepackage{multicol}
\usepackage{draftwatermark}

%==================================================
% COURSE METADATA (EDIT THESE FOR EACH MODULE)
%==================================================
\newcommand{\CourseCode}{QRS1001}
\newcommand{\CourseName}{Sample Course}
\newcommand{\DocType}{Revision Notes}
\newcommand{\AcademicYear}{Academic Year 2025--2026}

\title{
\Huge \CourseCode\ \CourseName \\[0.3em]
\Large \AcademicYear
}
\author{
  \textit{\href{https://qrsntu.org}{Quantitative Research Society @NTU}}
}
\date{\today}


%==================================================
% TCOLORBOX THEOREM STYLES
%==================================================

% ---------- DEFINITION: dark green, title in darker band ----------
\newtcbtheorem[number within=section]{definition}{Definition}%
{enhanced,
 breakable,
 colback=green!5!white,        % light interior
 colframe=green!40!black,      % border
 colbacktitle=green!50!black,  % dark title band
 coltitle=white,               % title text colour
 fonttitle=\bfseries,
 title filled,                 % puts title inside the band
 boxed title style={
   size=small,
   top=2pt, bottom=2pt,
   left=6pt, right=6pt,
   arc=1pt, outer arc=1pt
 },
 boxrule=0.9pt,
 arc=1pt,
 left=8pt, right=8pt, top=10pt, bottom=8pt,
 before skip=12pt, after skip=12pt
}{def}

% ---------- THEOREM: blue, title in darker band ----------
\newtcbtheorem[number within=section]{theorem}{Theorem}%
{enhanced,
 breakable,
 colback=blue!3!white,
 colframe=blue!40!black,
 colbacktitle=blue!55!black,
 coltitle=white,
 fonttitle=\bfseries,
 title filled,
 boxed title style={
   size=small,
   top=2pt, bottom=2pt,
   left=6pt, right=6pt,
   arc=1pt, outer arc=1pt
 },
 boxrule=0.9pt,
 arc=1pt,
 left=8pt, right=8pt, top=10pt, bottom=8pt,
 before skip=12pt, after skip=12pt
}{thm}

% ---------- ENVIRONMENTS USING TCOLORBOXENVIRONMENT ----------
\tcolorboxenvironment{identity}{
  enhanced,
  breakable,
  colback=green!5!white,
  colframe=green!60!black,
  fonttitle=\bfseries,
  boxrule=1pt,
  arc=3pt,
  left=8pt, right=8pt, top=8pt, bottom=8pt,
  before skip=12pt, after skip=12pt
}

\tcolorboxenvironment{principle}{
  enhanced,
  breakable,
  colback=red!5!white,
  colframe=red!60!black,
  fonttitle=\bfseries,
  boxrule=1pt,
  arc=3pt,
  left=8pt, right=8pt, top=8pt, bottom=8pt,
  before skip=12pt, after skip=12pt
}

%==================================================
% CUSTOM TCOLORBOX ENVIRONMENTS (WITH NAMES)
%==================================================
% Optional argument [1][] lets you override the default title if needed:
%   \begin{keyformula}[title={Key Formula: Demand Curve}] ... \end{keyformula}

\newtcolorbox{keyformula}[1][]{
    enhanced,
    breakable,
    colback=blue!5!white,
    colframe=blue!75!black,
    title=Key Formula,
    fonttitle=\bfseries,
    boxrule=0.9pt,
    arc=2pt,
    left=8pt, right=8pt, top=8pt, bottom=8pt,
    #1
}

% *** CHANGED TO BROWN AS REQUESTED ***
\newtcolorbox{examplebox}[1][]{
    enhanced,
    breakable,
    colback=brown!5!white,
    colframe=brown!75!black,
    title=Worked Example,
    fonttitle=\bfseries,
    boxrule=0.9pt,
    arc=2pt,
    left=8pt, right=8pt, top=8pt, bottom=8pt,
    #1
}

\newtcolorbox{commonmistake}[1][]{
    enhanced,
    breakable,
    colback=red!5!white,
    colframe=red!75!black,
    title=Common Mistake,
    fonttitle=\bfseries,
    boxrule=0.9pt,
    arc=2pt,
    left=8pt, right=8pt, top=8pt, bottom=8pt,
    #1
}

\newtcolorbox{policynote}[1][]{
    enhanced,
    breakable,
    colback=yellow!10!white,
    colframe=orange!75!black,
    title=Policy Application,
    fonttitle=\bfseries,
    boxrule=0.9pt,
    arc=2pt,
    left=8pt, right=8pt, top=8pt, bottom=8pt,
    #1
}

\newtcolorbox{examtip}[1][]{
    enhanced,
    breakable,
    colback=purple!5!white,
    colframe=purple!75!black,
    title=Exam Focus,
    fonttitle=\bfseries,
    boxrule=0.9pt,
    arc=2pt,
    left=8pt, right=8pt, top=8pt, bottom=8pt,
    #1
}

%==================================================
% HEADER / FOOTER / WATERMARK
%==================================================



\pagestyle{fancy}
\fancyhf{}
\fancyhead[L]{\CourseCode\ \CourseName{} -- \DocType}
\fancyhead[R]{\href{https://qrsntu.org}{QRS@NTU}}
\fancyfoot[C]{\thepage}
\renewcommand{\headrulewidth}{0.4pt}
\renewcommand{\footrulewidth}{0pt}

% Watermark
\SetWatermarkText{QRS@NTU — qrsntu.org}
\SetWatermarkScale{1}
\SetWatermarkLightness{0.99}

%------------- HEADER & FOOTER (for page 2 onward) -------------
\fancypagestyle{main}{
  \fancyhf{}
  \fancyhead[L]{\CourseCode\ \CourseName{} -- \DocType}
  \fancyhead[R]{\href{https://qrsntu.org}{QRS@NTU}}
  \fancyfoot[C]{\thepage}
  \renewcommand{\headrulewidth}{0.4pt}
  \renewcommand{\footrulewidth}{0pt}
}

%------------- SIMPLE FOOTER (page 1 only) -------------
\fancypagestyle{plainfirst}{
  \fancyhf{}
  \renewcommand{\headrulewidth}{0pt}
  \renewcommand{\footrulewidth}{0pt}
  \fancyfoot[C]{\scriptsize\textcolor{gray}{\textit{For academic reference only.
 This document is provided for educational purposes and does not constitute an 
 official question paper, answer key, or any formally authorised academic 
 material. Its content is not guaranteed to be complete or accurate.}}}
}

%------------- HYPERREF METADATA -------------
\hypersetup{
    colorlinks=false,
    linkcolor=blue,
    filecolor=magenta,
    urlcolor=blue,
    pdftitle={\CourseCode\ \CourseName{} -- \DocType},
    pdfauthor={Quantitative Research Society, NTU},
    pdfsubject={\CourseName},
    pdfkeywords={microeconomics, QRS@NTU, revision notes}
}

%==================================================
% DOCUMENT
%==================================================
\begin{document}
\maketitle
\thispagestyle{plainfirst}

\begin{abstract}
\noindent
Comprehensive revision notes for \CourseCode\ \CourseName, integrating theoretical concepts, key formulas, worked examples, and exam-focused applications.
\end{abstract}

\vspace{0.5cm}
\tableofcontents
\newpage

\newgeometry{left=0.7in,right=0.7in,top=1in,bottom=1in}
\fontsize{10.5pt}{11.5pt}\selectfont
\pagestyle{main}

%====================================================================
\section*{How to Use This Document}
\addcontentsline{toc}{section}{How to Use This Document}
%====================================================================

\subsection*{Structure of the Notes}

These notes are designed to be used together with your lectures and tutorials. Each chapter is organised into:

\begin{enumerate}
    \item \textbf{Core Theory} — formal definitions, results and economic intuition.
    \item \textbf{Mathematical Framework} — formulas, algebra and derivations.
    \item \textbf{Worked Examples} — tutorial-style questions with full solutions.
    \item \textbf{Applications} — policy or real-world interpretations where relevant.
    \item \textbf{Exam Focus} — high-yield concepts, typical traps and checklist items.
\end{enumerate}

A good way to revise:
\begin{itemize}
    \item First pass: read Core Theory + Key Formula boxes.
    \item Second pass: attempt Worked Examples yourself \emph{before} checking solutions.
    \item Final pass: review Exam Focus and Common Mistake boxes the week before the exam.
\end{itemize}

\subsection*{Colour and Box System (Legend)}

The document uses a colour-coded system. Each box has a \textbf{name} at the top so you can quickly recognise its purpose:

\begin{center}
\begin{tabular}{|l|l|p{7cm}|}
\hline
\textbf{Colour} & \textbf{Box Name} & \textbf{Use This For} \\
\hline
\colorbox{blue!5}{Blue}   & \textbf{Key Formula}      & Core equations and identities that you should memorise or be able to derive quickly. \\
\hline
\colorbox{brown!5}{Brown} & \textbf{Worked Example}   & Fully worked questions. Try on your own before reading the solution. \\
\hline
\colorbox{red!5}{Red}     & \textbf{Common Mistake}   & Typical errors, misleading shortcuts and conceptual traps to avoid in exams. \\
\hline
\colorbox{yellow!10}{Yellow} & \textbf{Policy Application} & Real-world or policy-related interpretations to help build intuition and essay points. \\
\hline
\colorbox{purple!5}{Purple} & \textbf{Exam Focus}      & High-yield checklists, ``must-know'' results and exam strategy tips. \\
\hline
\end{tabular}
\end{center}

\noindent
When you print in black and white, you can still identify boxes by their \textbf{title text}:
\begin{itemize}
    \item Look for the name at the top: \textit{Key Formula}, \textit{Worked Example}, \textit{Common Mistake}, \textit{Policy Application}, \textit{Exam Focus}.
    \item In digital version, use the colour + name; in printed version, rely mainly on the name.
\end{itemize}

\subsection*{Suggested Revision Flow}

\begin{examtip}
\begin{itemize}
    \item Use the \textbf{Key Formula} boxes to build a one-page cheat sheet.
    \item Re-do all \textbf{Worked Examples} at least once without looking at the solution.
    \item Read \textbf{Common Mistake} boxes immediately before attempting past-year papers.
    \item Use \textbf{Exam Focus} boxes as your final 1--2 day cram checklist.
\end{itemize}
\end{examtip}

\newpage
%==================================================
\part{Sample Part Title}
%==================================================

\section{Sample Topic}

This section shows how to use the different box environments in practice.

\begin{definition}{Sample Definition}{sample-definition}
A \emph{sample object} is anything you use here to test the formatting of this template. Replace this with your actual course content.
\end{definition}

\begin{theorem}{Sample Theorem}{sample-theorem}
For every sample object in these notes, there exists at least one student who will replace it with real content.
\end{theorem}

\begin{keyformula}
For a linear demand curve of the form
\[
Q = a - bP,
\]
the corresponding inverse demand and total revenue are
\[
P = \frac{a - Q}{b}, \qquad TR(Q) = P \cdot Q.
\]
\end{keyformula}

\begin{examplebox}
\textbf{Worked Example: Replace With Real Question}

This is a placeholder example showing the brown \emph{Worked Example} box.

\begin{enumerate}
    \item State your question clearly.
    \item Show all working step by step.
    \item Box your final numerical or algebraic answer.
\end{enumerate}
\end{examplebox}

\begin{commonmistake}
Students often skip intermediate algebra steps because they ``can do it in their head'' and then lose marks due to sign errors or missing terms. In exam conditions, write at least one clear intermediate line for each non-trivial step.
\end{commonmistake}

\begin{policynote}
When you see a result in microeconomics, always ask: \emph{``What real-world policy or market does this apply to?''} Linking theory to context makes it easier to recall under pressure.
\end{policynote}

\begin{examtip}
Before the exam, make a list of:
\begin{itemize}
    \item Top 10 formulas you must know.
    \item 3--5 ``classic'' Common Mistakes you personally tend to make.
    \item 1--2 favourite examples you can write quickly if an open-ended question appears.
\end{itemize}
\end{examtip}

\section{Hamilton--Jacobi--Bellman (HJB) Equation}
\label{sec:hjb}
%==========================================================

\subsection{Stochastic Control Setup}

We consider a controlled diffusion process $(X_t)_{t\ge 0}$ satisfying
\[
  dX_t = b(X_t,u_t)\,dt + \sigma(X_t,u_t)\,dW_t,
\]
where
\begin{itemize}
  \item $X_t$ is the \textbf{state} (e.g.\ wealth, capital, portfolio value),
  \item $u_t \in \mathcal{U}$ is the \textbf{control} (e.g.\ portfolio weight, consumption rate),
  \item $W_t$ is a standard Brownian motion.
\end{itemize}

Given a control policy $(u_t)_{t\ge 0}$, define the discounted performance criterion
\[
  J^{u}(x) = \mathbb{E}_x\Bigg[\int_0^\infty e^{-\rho t} f(X_t,u_t)\,dt\Bigg],
\]
where $\rho>0$ is a discount rate and $f$ is the running reward (or utility) function.

\begin{definition}{Value function}{hjb_value}
For an initial state $x$, the \textbf{value function} is
\[
  V(x) := \sup_{u \in \mathcal{A}} J^{u}(x),
\]
where $\mathcal{A}$ is the set of admissible control policies. It represents
the maximal expected discounted reward achievable from state $x$.
\end{definition}

\subsection{HJB Equation (Informal Statement)}

\begin{theorem}{Hamilton--Jacobi--Bellman Equation (Informal)}{hjb_informal}
Suppose the control problem above is well behaved and $V$ is sufficiently smooth.
Then $V$ satisfies the \textbf{stationary HJB equation}
\[
  \rho V(x)
  = \sup_{u\in\mathcal{U}}
    \Big\{ f(x,u)
          + V_x(x)\,b(x,u)
          + \tfrac12 V_{xx}(x)\,\sigma^2(x,u)
    \Big\},
\]
where $V_x$ and $V_{xx}$ denote the first and second derivatives of $V$.
\end{theorem}

\begin{keyformula}
\textbf{Stationary HJB (1D, infinite horizon)}

\[
  \rho V(x)
  = \sup_{u\in\mathcal{U}}
    \left[
      f(x,u)
      + V_x(x)\,b(x,u)
      + \frac{1}{2} V_{xx}(x)\,\sigma^2(x,u)
    \right].
\]

\begin{itemize}
  \item LHS: \(\rho V(x)\) = ``required return'' on value.
  \item RHS: instantaneous payoff \(f\) + expected capital gain
        via drift and diffusion.
  \item Optimal control \(u^\ast(x)\) maximizes the bracket term pointwise in \(u\).
\end{itemize}
\end{keyformula}

\begin{examtip}
When writing down the HJB:
\begin{enumerate}
  \item Write the SDE for \(X_t\) clearly.
  \item Define the objective functional \(J^{u}(x)\).
  \item State the value function \(V(x)\).
  \item Then write the HJB as
        \(\rho V = \sup_{u}\{\text{running reward} + \text{generator of }X\text{ applied to }V\}\).
\end{enumerate}
Marks are usually given for the structure even if you do not solve the PDE.
\end{examtip}

\subsection{Simple Deterministic HJB Example}

To see the mechanics without stochastic noise, consider a \emph{deterministic}
control problem:
\[
  \dot{x}(t) = u(t), \qquad x(0) = x,
\]
with cost functional
\[
  J^{u}(x) = \int_0^\infty e^{-\rho t}
  \left( -\tfrac12 u(t)^2 \right)\,dt,
\]
where $\rho>0$ and $u(t)\in\mathbb{R}$.

Here,
\[
  f(x,u) = -\tfrac12 u^2, \qquad b(x,u)=u, \qquad \sigma \equiv 0.
\]

\begin{examplebox}
\textbf{Deterministic LQ Control: HJB and Optimal Control}

\begin{enumerate}[label=(\alph*)]
  \item \textbf{HJB equation.}  
  The HJB becomes
  \[
    \rho V(x)
    = \sup_{u\in\mathbb{R}}
      \left\{
        -\tfrac12 u^2 + V_x(x)\,u
      \right\}.
  \]

  \item \textbf{Optimize over \(u\).}  
  For fixed \(x\), we maximize
  \[
    g(u) = -\tfrac12 u^2 + V_x(x)\,u.
  \]
  First-order condition:
  \[
    g'(u) = -u + V_x(x) = 0
    \quad\Rightarrow\quad
    u^\ast(x) = V_x(x).
  \]
  Second derivative \(g''(u)=-1<0\) confirms a maximum.

  \item \textbf{Plug back into HJB.}  
  Substitute \(u^\ast(x)\) into the HJB:
  \[
    \rho V(x)
    = -\tfrac12 (V_x(x))^2 + V_x(x)\,V_x(x)
    = \tfrac12 \bigl(V_x(x)\bigr)^2.
  \]
  So \(V\) must satisfy the first-order ODE
  \[
    \rho V(x) = \frac{1}{2}\bigl(V_x(x)\bigr)^2.
  \]

  \item \textbf{Takeaway.}  
  Even in this simple example, the HJB:
  \begin{itemize}
    \item comes from \(\rho V = \sup_u\{f + V_x b\}\),
    \item yields an optimal feedback control \(u^\ast(x)\),
    \item and reduces the dynamic optimization problem to solving a nonlinear ODE for \(V\).
  \end{itemize}
\end{enumerate}
\end{examplebox}

\begin{commonmistake}
In exam settings, a very common mistake is to forget the discount term
\(\rho V(x)\) on the left-hand side, or to accidentally write
\(\rho V'(x)\) instead of \(\rho V(x)\). Always check:
\[
\boxed{\text{LHS is } \rho V, \text{ not } \rho V'.}
\]
\end{commonmistake}



%==========================================================
% (You can continue with more sections for your actual course here.)
%==========================================================
%====================================================
\appendix
%====================================================

\section{Examination Preparation Guide}

\subsection{Overview of Practice Resources}

The Quantitative Research Society has compiled a focused preparation pack for
the HJB / stochastic control component of the course. The aim is to give you
\emph{structured, exam-style} practice on all major skills:

\begin{center}
\begin{tabular}{l|c|c|p{5.5cm}}
\toprule
\textbf{Resource} & \textbf{Type} & \textbf{Length} & \textbf{Description} \\
\midrule
Problem Set A & Short questions & 15--20 Qs & Concept checks on dynamic programming, value functions, Bellman principle \\
Problem Set B & Calculations & 10--15 Qs & HJB derivations, simple SDEs, generators, verification arguments \\
Problem Set C & Applications & 8--10 Qs & Portfolio / consumption problems, deterministic control, LQ examples \\
Mock Exam 1 & Full paper & 90--120 min & Mixed conceptual and computational; similar to a typical final exam section \\
Mock Exam 2 & Full paper & 90--120 min & Slightly more theoretical; emphasises proofs and derivations of HJB \\
\bottomrule
\end{tabular}
\end{center}

\subsection{Suggested Study Sequence}

\begin{enumerate}[noitemsep]
  \item \textbf{Read}: Section~\ref{sec:hjb} of these notes once quickly to see the “big picture”.
  \item \textbf{Drill}: Work through Problem Sets A and B with solutions closed.
  \item \textbf{Consolidate}: Attempt Problem Set C to connect formulas to economic/financial intuition.
  \item \textbf{Simulate exam}: Sit Mock Exam 1 under timed conditions.
  \item \textbf{Refine}: Review your mistakes; then sit Mock Exam 2 for final tuning.
\end{enumerate}

\begin{examtip}
In many exams, the HJB part is \emph{high-yield but template-based}. You typically
receive marks for:
\begin{enumerate}[noitemsep]
  \item Stating the control problem clearly (state, control, objective).
  \item Writing down the correct HJB equation with all terms present.
  \item Identifying the optimal control via a first-order condition or maximisation step.
  \item Simplifying the resulting ODE/PDE and, where required, solving it.
\end{enumerate}
Even if you cannot solve the equation completely, a \emph{correct HJB setup}
usually secures a substantial portion of the marks.
\end{examtip}

%====================================================

%====================================================

\begingroup
\newgeometry{left=0.7in,right=0.7in,top=1in,bottom=1in}
\fontsize{10pt}{11pt}\selectfont
\section{Formula \& Key Results Reference}
\begin{multicols}{2}

%--------------------------------------------
\subsection*{Dynamic Programming and Value Functions}
%--------------------------------------------

\noindent\textbf{Value function (infinite horizon, discounted):}
\[
V(x) = \sup_{u\in\mathcal{A}}
\mathbb{E}_x\!\left[\int_0^\infty e^{-\rho t} f(X_t,u_t)\,dt\right].
\]

\noindent\textbf{Dynamic programming principle (informal):}
\[
V(x) = \sup_{u}
\mathbb{E}_x\!\left[
  \int_0^h e^{-\rho t} f(X_t,u_t)\,dt
  + e^{-\rho h} V(X_h)
\right]\!.
\]

\noindent\textbf{Bellman recursion (discrete time):}
\[
V(x) = \sup_{u\in U(x)} \Big\{
  r(x,u) + \beta\,\mathbb{E}\big[V(X' ) \mid x,u\big]
\Big\}.
\]

%--------------------------------------------
\subsection*{Controlled SDE and Generator}
%--------------------------------------------

\noindent\textbf{Controlled SDE (1D):}
\[
dX_t = b(X_t,u_t)\,dt + \sigma(X_t,u_t)\,dW_t.
\]

\noindent\textbf{Generator of $(X_t)$ under control $u$:}
\[
\mathcal{L}^u \phi(x)
= b(x,u)\,\phi_x(x)
+ \frac12 \sigma^2(x,u)\,\phi_{xx}(x).
\]

\noindent\textbf{Key identity for HJB:}
\[
\mathcal{L}^u V(x)
= V_x(x)\,b(x,u)
+ \tfrac12 V_{xx}(x)\,\sigma^2(x,u).
\]

%--------------------------------------------
\subsection*{HJB Equations}
%--------------------------------------------

\noindent\textbf{Stationary (infinite-horizon) HJB:}
\[
\rho V(x) =
\sup_{u\in\mathcal{U}}
\left\{
  f(x,u) + \mathcal{L}^u V(x)
\right\}.
\]

\noindent\textbf{Finite-horizon HJB (backward in $t$):}
\[
\begin{cases}
V_t(t,x)
+ \displaystyle \sup_{u\in\mathcal{U}}
\bigl\{
  f(t,x,u) + \mathcal{L}^u V(t,x)
\bigr\}
= 0,\\[0.4em]
V(T,x) = g(x) \quad \text{(terminal condition)}.
\end{cases}
\]

\noindent\textbf{Deterministic control (no noise):}
If $\dot{x} = b(x,u)$ and cost rate $f(x,u)$, then
\[
V_t(t,x)
+ \sup_{u}
\bigl\{ f(x,u) + V_x(t,x)\,b(x,u)\bigr\}
= 0.
\]

%--------------------------------------------
\subsection*{It\^o’s Lemma (Scalar Version)}
%--------------------------------------------

\noindent If
\[
dX_t = \mu(t,X_t)\,dt + \sigma(t,X_t)\,dW_t,
\]
and $Y_t = \phi(t,X_t)$, then
\[
dY_t
= \left(
    \phi_t
    + \mu\,\phi_x
    + \tfrac12 \sigma^2\,\phi_{xx}
  \right)\!dt
+ \sigma\,\phi_x\,dW_t.
\]

\noindent\textbf{Shortcut memory:}
\[
\boxed{
d\phi = \phi_t\,dt + \phi_x\,dX + \tfrac12 \phi_{xx}\,(dX)^2
}
\]
and for an SDE, $(dW_t)^2 = dt$, all higher powers vanish.

%--------------------------------------------
\subsection*{Verification Theorem (Very Informal)}
%--------------------------------------------

\noindent\textbf{Template:}
\begin{itemize}[noitemsep,leftmargin=*]
  \item Guess a smooth candidate $V$.
  \item Check that $V$ satisfies the HJB equation.
  \item Show that the control $u^\ast(x)$ that attains the supremum in the HJB
        is admissible.
  \item Conclude: $V$ is the true value function and $u^\ast$ is optimal.
\end{itemize}

\noindent\textbf{Sign check:} For a \emph{maximum} problem, the HJB uses
\(\sup\); for a \emph{minimum} problem, replace \(\sup\) by \(\inf\) and
reverse inequalities accordingly.

%--------------------------------------------
\subsection*{Deterministic LQ-Control (Memory Hook)}
%--------------------------------------------

\noindent System:
\[
\dot{x}(t) = a x(t) + b u(t),
\quad
J(u) = \int_0^\infty e^{-\rho t}
\bigl( q x(t)^2 + r u(t)^2 \bigr)\,dt.
\]

\noindent Ansatz:
\[
V(x) = K x^2
\quad\Rightarrow\quad
u^\ast(x) = -\frac{b}{r} V_x(x) = -\frac{2Kb}{r} x.
\]

\noindent HJB reduces to an algebraic Riccati equation for $K$:
\[
\rho K x^2
=
q x^2 +
(2Kx)(a x + b u^\ast)
+ r \bigl(u^\ast\bigr)^2.
\]

\noindent\textbf{Key idea:} Quadratic cost $\Rightarrow$ quadratic value function
$\Rightarrow$ \emph{linear} optimal feedback rule.

%--------------------------------------------
\subsection*{Concept Checklist (Typical Exam Questions)}
%--------------------------------------------

\begin{itemize}[noitemsep,leftmargin=*]
  \item Can you \textbf{state} the control problem precisely
        (state, control, objective, horizon)?
  \item Can you write down the \textbf{generator} $\mathcal{L}^u$ of a given SDE?
  \item Can you derive the correct \textbf{HJB} (finite or infinite horizon)?
  \item Can you find $u^\ast(x)$ by \textbf{maximizing the Hamiltonian} /
        the inner bracket in the HJB?
  \item Can you explain in words what HJB / Bellman principle means
        (``optimal policy is myopic with respect to $V$'')?
  \item In deterministic problems, can you connect HJB and the
        \textbf{Pontryagin maximum principle} (same spirit, different language)?
\end{itemize}

\end{multicols}

\restoregeometry
\endgroup


\end{document}
